\allowdisplaybreaks
\begin{center}
	\hrule
	\vspace{.4cm}
 \begin{tabular*}{\textwidth}{@{}l@{}|@{\extracolsep{0.6in}}r@{}}%
\parbox{4.25in}{\raggedright{\includegraphics[width=.9\linewidth]{ictp-pwf.pdf}}} &
\parbox[c][]{4in}{{\Large\textbf{Md Akiful Islam Zawad} \par}
                    { Brac University \par}
                    { KHA 224, Progati Sarani, Merul Badda, \par}
                    { Dhaka 1212, Bangladesh \par}
                    { \href{ext.akiful.islam@bracu.ac.bd}{ext.akiful.islam@bracu.ac.bd}} \par}
\end{tabular*}\vspace{.3in}
\definecolor{ceruleanblue}{rgb}{0.16, 0.32, 0.6}
	\LARGE\scshape\textbf{\textcolor{ceruleanblue}{Physics For Bangladesh: School On Quantum Field Theory}}
\end{center}
\hrule\vspace{.25in}
{\large\textbf{Problem Sheet:}\ \textsc{1} \hspace{\hfill} \large\textbf{Due Date:} \today\\
	\hline\hline
 %----------------------------
\paragraph*{Problem Set 4} %\hfill \newline
\\
The Fourier decompositions of a real scalar field and its conjugate momentum are (note that they are at $t = 0$, as compared to the lectures)
    \begin{align*}
        \phi(\vec{x}) &= \bigints \frac{d^3p}{(2\pi)^3} \frac{1}{\sqrt{2E_{\mathbf{p}}}} \left[a_{\mathbf{p}} e^{i\mathbf{p} \cdot \vec{x}} + a^\dagger_{\mathbf{p}} e^{-i\mathbf{p} \cdot \vec{x}} \right]\\
        \pi(\vec{x}) &= \bigints \frac{d^3p}{(2\pi)^3} (-i)\sqrt{\frac{E_{\mathbf{p}}}{2}} \left[a_{\mathbf{p}} e^{i\mathbf{p} \cdot \vec{x}} - a^\dagger_{\mathbf{p}} e^{-i\mathbf{p} \cdot \vec{x}} \right]
    \end{align*}
    with (equal time at $t = 0$) commutation relationships
    \begin{align*}
        [\phi(\vec{x}), \phi(\vec{y})] = [\pi(\vec{x}), \pi(\vec{y})] = 0 \quad \text{and} \quad [\phi(\vec{x}), \pi(\vec{y})] = i\delta^{(3)}(\vec{x} - \vec{y})
    \end{align*}
    implying
    \begin{align*}
        [a_{\mathbf{p}}, a_{\vec{q}}] = [a^\dagger_{\mathbf{p}}, a^\dagger_{\vec{q}}] = 0 \quad \text{and} \quad [a_{\mathbf{p}}, a^\dagger_{\vec{q}}] = (2\pi)^3 \delta^{(3)}(\mathbf{p} - \vec{q})
    \end{align*}
    \begin{enumerate}
        \item Generalize above analysis for the complex scalar field whose Lagrangian is given in problem 2.
        \item Compute the Feynman propagator for the complex scalar field defined by $\langle 0|T(\Phi(x) \Phi^\dagger(y))|0 \rangle$.
        \item Use the mode expansion of the field and derive the propagator in momentum space.
        \item Lastly, compute the time-ordered product $T(\Phi(x) \Phi^\dagger(y) \Phi(z))$.
    \end{enumerate}
\bigskip\bigskip\hline\hline\bigskip
\section*{Solution}
\subsection*{(a) Generalization to Complex Scalar Field}
Consider a complex scalar field \(\Phi(\vec{x})\) with the Lagrangian density:
\begin{align}
    \mathcal{L} = \partial_\mu \Phi^* \partial^\mu \Phi - m^2 \Phi^* \Phi - \frac{\lambda}{2} (\Phi^* \Phi)^2.
\end{align}
For simplicity, we will first consider the free case, where \(\lambda = 0\):
\begin{align}
    \mathcal{L}_{\text{free}} = \partial_\mu \Phi^* \partial^\mu \Phi - m^2 \Phi^* \Phi.
\end{align}
The Euler-Lagrange equation for the field \(\Phi\) is given by:
\begin{align}
    \frac{\partial \mathcal{L}}{\partial \Phi^*} - \partial_\mu \left( \frac{\partial \mathcal{L}}{\partial (\partial_\mu \Phi^*)} \right) = 0.
\end{align}
For \(\Phi\), we have:
\begin{align}
    \frac{\partial \mathcal{L}}{\partial \Phi^*} &= -m^2 \Phi, \\
    \frac{\partial \mathcal{L}}{\partial (\partial_\mu \Phi^*)} &= \partial^\mu \Phi.
\end{align}
Thus,
\begin{align}
    \partial_\mu \left( \frac{\partial \mathcal{L}}{\partial (\partial_\mu \Phi^*)} \right) &= \partial_\mu \left( \partial^\mu \Phi \right) = \Box \Phi,
\end{align}
where \(\Box = \partial_\mu \partial^\mu\) is the d'Alembertian operator. The equation of motion is:
\begin{align}
    (\Box + m^2) \Phi = 0.
\end{align}
The general solution to the Klein-Gordon equation can be expressed as a superposition of plane waves. We write the mode expansion for the field \(\Phi(\vec{x})\) as:
\begin{align}
    \Phi(\vec{x}) &= \bigints\frac{d^3p}{(2\pi)^3} \frac{1}{\sqrt{2E_{\mathbf{p}}}} \left( a_{\mathbf{p}} e^{i \mathbf{p} \cdot \vec{x}} + b^\dagger_{\mathbf{p}} e^{-i \mathbf{p} \cdot \vec{x}} \right),
\end{align}
where \(E_{\mathbf{p}} = \sqrt{\mathbf{p}^2 + m^2}\) is the energy of the particle.

The conjugate momentum \(\pi(\vec{x})\) is defined as:
\begin{align}
    \pi(\vec{x}) = \frac{\partial \mathcal{L}}{\partial (\partial_0 \Phi^*)}.
\end{align}
For our Lagrangian density, we get:
\begin{align}
    \pi(\vec{x}) &= \frac{\partial}{\partial t} \Phi^*(\vec{x}).
\end{align}
To find \(\pi(\vec{x})\), we first need the mode expansion for \(\Phi^*(\vec{x})\). The complex conjugate of \(\Phi(\vec{x})\) is:
\begin{align}
    \Phi^*(\vec{x}) &= \bigints\frac{d^3p}{(2\pi)^3} \frac{1}{\sqrt{2E_{\mathbf{p}}}} \left( a^\dagger_{\mathbf{p}} e^{-i \mathbf{p} \cdot \vec{x}} + b_{\mathbf{p}} e^{i \mathbf{p} \cdot \vec{x}} \right).
\end{align}
Taking the time derivative:
\begin{align}
    \pi(\vec{x}) &= \frac{\partial}{\partial t} \Phi^*(\vec{x}) \\
    &= \bigints\frac{d^3p}{(2\pi)^3} \frac{1}{\sqrt{2E_{\mathbf{p}}}} \left( -i E_{\mathbf{p}} a^\dagger_{\mathbf{p}} e^{-i \mathbf{p} \cdot \vec{x}} + i E_{\mathbf{p}} b_{\mathbf{p}} e^{i \mathbf{p} \cdot \vec{x}} \right) \\
    &= \bigints\frac{d^3p}{(2\pi)^3} \left(-i \sqrt{\frac{E_{\mathbf{p}}}{2}} \left( a^\dagger_{\mathbf{p}} e^{-i \mathbf{p} \cdot \vec{x}} - b_{\mathbf{p}} e^{i \mathbf{p} \cdot \vec{x}} \right)\right).
\end{align}
The mode expansions for the field \(\Phi(\vec{x})\) and its conjugate momentum \(\pi(\vec{x})\) are:
\begin{align}
    \Phi(\vec{x}) &= \bigints\frac{d^3p}{(2\pi)^3} \frac{1}{\sqrt{2E_{\mathbf{p}}}} \left( a_{\mathbf{p}} e^{i \mathbf{p} \cdot \vec{x}} + b^\dagger_{\mathbf{p}} e^{-i \mathbf{p} \cdot \vec{x}} \right), \\
    \pi(\vec{x}) &= \bigints\frac{d^3p}{(2\pi)^3} \left(-i \sqrt{\frac{E_{\mathbf{p}}}{2}} \left( a^\dagger_{\mathbf{p}} e^{-i \mathbf{p} \cdot \vec{x}} - b_{\mathbf{p}} e^{i \mathbf{p} \cdot \vec{x}} \right)\right).
\end{align}
We now want to check the commutation relation:
\begin{align}
    [\phi(\vec{x}), \phi(\vec{y})] &= [\pi(\vec{x}), \pi(\vec{y})].
\end{align}
\begin{align*}
    [\phi(\vec{x}), \phi(\vec{y})] &= \left[ \bigints\frac{d^3p}{(2\pi)^3} \frac{1}{\sqrt{2E_{\mathbf{p}}}} \left(a_{\mathbf{p}} e^{i \mathbf{p} \cdot \vec{x}} + b^\dagger_{\mathbf{p}} e^{-i \mathbf{p} \cdot \vec{x}}\right) \bigints\frac{d^3p^\prime}{(2\pi)^3} \frac{1}{\sqrt{2E_{\mathbf{p}^\prime}}} \left(a_{\mathbf{p}^\prime} e^{i \mathbf{p}^\prime \cdot \vec{y}} + b^\dagger_{\mathbf{p}^\prime} e^{-i \mathbf{p}^\prime \cdot \vec{y}}\right) \right] \\
    &= \bigints\frac{d^3p}{(2\pi)^3} \frac{d^3p^\prime}{(2\pi)^3}\left[\frac{1}{\sqrt{2E_{\mathbf{p}}}} \frac{1}{\sqrt{2E_{\mathbf{p}^\prime}}} \left[a_{\mathbf{p}}, a_{\mathbf{p}^\prime} \right] e^{i \mathbf{p} \cdot \vec{x}} e^{i \mathbf{p}^\prime \cdot \vec{y}}+ \frac{1}{\sqrt{2E_{\mathbf{p}}}} \frac{1}{\sqrt{2E_{\mathbf{p}^\prime}}} \left[a_{\mathbf{p}}, b^\dagger_{\mathbf{p}^\prime} \right] e^{i \mathbf{p} \cdot \vec{x}} e^{-i \mathbf{p}^\prime \cdot \vec{y}}\right] \\
    & \quad + \bigints\frac{d^3p}{(2\pi)^3} \frac{d^3p^\prime}{(2\pi)^3}\left[\frac{1}{\sqrt{2E_{\mathbf{p}}}} \frac{1}{\sqrt{2E_{\mathbf{p}^\prime}}} \left[b^\dagger_{\mathbf{p}}, a_{\mathbf{p}^\prime} \right] e^{-i \mathbf{p} \cdot \vec{x}} e^{i \mathbf{p}^\prime \cdot \vec{y}} + \frac{1}{\sqrt{2E_{\mathbf{p}}}} \frac{1}{\sqrt{2E_{\mathbf{p}^\prime}}} \left[b^\dagger_{\mathbf{p}}, b^\dagger_{\mathbf{p}^\prime} \right] e^{-i \mathbf{p} \cdot \vec{x}} e^{-i \mathbf{p}^\prime \cdot \vec{y}} \right] \\
    &= 0.
\end{align*}
\begin{align*}
    [\pi(\vec{x}), \pi(\vec{y})] &= \left[ \bigints\frac{d^3p}{(2\pi)^3} \left(-i \sqrt{\frac{E_{\mathbf{p}}}{2}} \left(a^\dagger_{\mathbf{p}} e^{-i \mathbf{p} \cdot \vec{x}} - b_{\mathbf{p}} e^{i \mathbf{p} \cdot \vec{x}}\right)\right)\bigints\frac{d^3p^\prime}{(2\pi)^3} \left(-i \sqrt{\frac{E_{\mathbf{p}^\prime}}{2}} \left(a^\dagger_{\mathbf{p}^\prime} e^{-i \mathbf{p}^\prime \cdot \vec{y}} - b_{\mathbf{p}^\prime} e^{i \mathbf{p}^\prime \cdot \vec{y}}\right)\right) \right] \\
    &= \bigints\frac{d^3p}{(2\pi)^3} \frac{d^3p^\prime}{(2\pi)^3} \left[\left(-i \sqrt{\frac{E_{\mathbf{p}}}{2}} \sqrt{\frac{E_{\mathbf{p}^\prime}}{2}}\right) \left[a^\dagger_{\mathbf{p}}, a_{\mathbf{p}^\prime}\right] e^{-i \mathbf{p} \cdot \vec{x}} e^{-i \mathbf{p}^\prime \cdot \vec{y}} + \left(-i \sqrt{\frac{E_{\mathbf{p}}}{2}} \sqrt{\frac{E_{\mathbf{p}^\prime}}{2}}\right) \left[a^\dagger_{\mathbf{p}}, b_{\mathbf{p}^\prime}\right] e^{-i \mathbf{p} \cdot \vec{x}} e^{i \mathbf{p}^\prime \cdot \vec{y}}\right] \\
    & \quad + \bigints\frac{d^3p}{(2\pi)^3} \frac{d^3p^\prime}{(2\pi)^3} \left[\left(-i \sqrt{\frac{E_{\mathbf{p}}}{2}} \sqrt{\frac{E_{\mathbf{p}^\prime}}{2}}\right) \left[b_{\mathbf{p}}, a^\dagger_{\mathbf{p}^\prime}\right] e^{i \mathbf{p} \cdot \vec{x}} e^{-i \mathbf{p}^\prime \cdot \vec{y}} + \left(-i \sqrt{\frac{E_{\mathbf{p}}}{2}} \sqrt{\frac{E_{\mathbf{p}^\prime}}{2}}\right) \left[b_{\mathbf{p}}, b_{\mathbf{p}^\prime}\right] e^{i \mathbf{p} \cdot \vec{x}} e^{i \mathbf{p}^\prime \cdot \vec{y}}\right] \\
    &= 0.
\end{align*}
with commutation relationships
    \begin{align*}
        [\phi(\vec{x}), \phi(\vec{y})] = [\pi(\vec{x}), \pi(\vec{y})] = 0 \quad \text{and} \quad [\phi(\vec{x}), \pi(\vec{y})] = i\delta^{(3)}(\vec{x} - \vec{y})
    \end{align*}
    implying
    \begin{align*}
        [a_{\mathbf{p}}, a_{\vec{q}}] = [a^\dagger_{\mathbf{p}}, a^\dagger_{\vec{q}}] = 0 \quad \text{and} \quad [a_{\mathbf{p}}, a^\dagger_{\vec{q}}] = (2\pi)^3 \delta^{(3)}(\mathbf{p} - \vec{q})
    \end{align*}
\subsection*{(b) Feynman Propagator for the Complex Scalar Field}
The Feynman propagator is defined as:
\begin{align}
    \Delta_F(x - y) &= \langle 0 | T(\Phi(x) \Phi^\dagger(y)) | 0 \rangle.
\end{align}
To compute this, we use the field expansions:
\begin{align}
    \Phi(x) &= \bigints \frac{d^3p}{(2\pi)^3} \frac{1}{\sqrt{2E_{\mathbf{p}}}} \left[ a_{\mathbf{p}} e^{i \mathbf{p} \cdot \vec{x}} + b^\dagger_{\mathbf{p}} e^{-i \mathbf{p} \cdot \vec{x}} \right], \\
    \Phi^\dagger(y) &= \bigints \frac{d^3p^\prime}{(2\pi)^3} \frac{1}{\sqrt{2E_{\mathbf{p}^\prime}}} \left[ a^\dagger_{\mathbf{p}^\prime} e^{-i \mathbf{p}^\prime \cdot \vec{y}} + b_{\mathbf{p}^\prime} e^{i \mathbf{p}^\prime \cdot \vec{y}} \right].
\end{align}

We then compute the vacuum expectation value:

\begin{align}
    \langle 0 | \Phi(x) \Phi^\dagger(y) | 0 \rangle &= \bigints \frac{d^3p}{(2\pi)^3} \frac{d^3p^\prime}{(2\pi)^3} \frac{1}{\sqrt{2E_{\mathbf{p}}}} \frac{1}{\sqrt{2E_{\mathbf{p}^\prime}}} \langle 0 | \left[ a_{\mathbf{p}} e^{i \mathbf{p} \cdot \vec{x}} + b^\dagger_{\mathbf{p}} e^{-i \mathbf{p} \cdot \vec{x}} \right] \left[ a^\dagger_{\mathbf{p}^\prime} e^{-i \mathbf{p}^\prime \cdot \vec{y}} + b_{\mathbf{p}^\prime} e^{i \mathbf{p}^\prime \cdot \vec{y}} \right] | 0 \rangle\notag\\
    &= \bigints \frac{d^3p}{(2\pi)^3} \frac{d^3p^\prime}{(2\pi)^3} \frac{1}{\sqrt{2E_{\mathbf{p}}}} \frac{1}{\sqrt{2E_{\mathbf{p}^\prime}}} \Bigg[ \langle 0 | a_{\mathbf{p}} a^\dagger_{\mathbf{p}^\prime} | 0 \rangle e^{i \mathbf{p} \cdot \vec{x}} e^{-i \mathbf{p}^\prime \cdot \vec{y}} + \langle 0 | a_{\mathbf{p}} b_{\mathbf{p}^\prime} e^{i \mathbf{p} \cdot \vec{x}} e^{i \mathbf{p}^\prime \cdot \vec{y}} | 0 \rangle \notag\\
    &\quad + \langle 0 | b^\dagger_{\mathbf{p}} a^\dagger_{\mathbf{p}^\prime} e^{-i \mathbf{p} \cdot \vec{x}} e^{-i \mathbf{p}^\prime \cdot \vec{y}} | 0 \rangle + \langle 0 | b^\dagger_{\mathbf{p}} b_{\mathbf{p}^\prime} e^{-i \mathbf{p} \cdot \vec{x}} e^{i \mathbf{p}^\prime \cdot \vec{y}} | 0 \rangle \Bigg].
\end{align}
Since \(\langle 0 | a_{\mathbf{p}} b_{\mathbf{p}^\prime} | 0 \rangle = \langle 0 | b^\dagger_{\mathbf{p}} a_{\mathbf{p}^\prime}^\dagger | 0 \rangle = 0\), the above equation simplifies to:
\begin{align}
    \langle 0 | \Phi(x) \Phi^\dagger(y) | 0 \rangle &= \bigints \frac{d^3p}{(2\pi)^3} \frac{d^3p^\prime}{(2\pi)^3} \frac{1}{\sqrt{2E_{\mathbf{p}}}} \frac{1}{\sqrt{2E_{\mathbf{p}^\prime}}} \Bigg[ \delta^3(\mathbf{p} - \mathbf{p}^\prime) e^{i \mathbf{p} \cdot \vec{x}} e^{-i \mathbf{p}^\prime \cdot \vec{y}} + \delta^3(\mathbf{p} - \mathbf{p}^\prime) e^{-i \mathbf{p} \cdot \vec{x}} e^{i \mathbf{p}^\prime \cdot \vec{y}} \Bigg]\notag\\
    &= \bigints \frac{d^3p}{(2\pi)^3} \frac{1}{2E_{\mathbf{p}}} \left[ e^{i \mathbf{p} \cdot (\vec{x} - \vec{y})} + e^{-i \mathbf{p} \cdot (\vec{x} - \vec{y})} \right].
\end{align}
Thus, the Feynman propagator is:
\begin{align}
    \Delta_F(x - y) &= \bigints \frac{d^3p}{(2\pi)^3} \frac{1}{2E_{\mathbf{p}}} e^{-i \mathbf{p} \cdot (\vec{x} - \vec{y})}.
\end{align}
\subsection*{(c) Feynman Propagator in momentum space}
The Feynman propagator (or two-point function) is:
\begin{align}
    \langle 0 | T\{\phi(x) \phi^\dagger(y)\} | 0 \rangle.
\end{align}
Expanding the fields using the mode expansions, we get:
\begin{align*}
    \langle 0 | T\{\phi(x) \phi^\dagger(y)\} | 0 \rangle &= \langle 0 \left| T \left\{ \left( \bigints\frac{d^3\mathbf{k}}{(2\pi)^{3/2}} \frac{1}{\sqrt{2E_k}} \left( a_{\mathbf{k}} e^{-ik \cdot x} + b^\dagger_{\mathbf{k}} e^{ik \cdot x} \right) \right) \left( \bigints\frac{d^3\mathbf{k}'}{(2\pi)^{3/2}} \frac{1}{\sqrt{2E_{k'}}} \left( b_{\mathbf{k}'} e^{-ik' \cdot y} + a^\dagger_{\mathbf{k}'} e^{ik' \cdot y} \right) \right) \right\} \right| 0 \rangle \\
    &= \bigints\frac{d^3\mathbf{k}}{(2\pi)^{3/2}} \frac{1}{\sqrt{2E_k}} \bigints\frac{d^3\mathbf{k}'}{(2\pi)^{3/2}} \frac{1}{\sqrt{2E_{k'}}} \left\langle 0 \left| \left(a_{\mathbf{k}} e^{-ik \cdot x} + b^\dagger_{\mathbf{k}} e^{ik \cdot x} \right) \left(b_{\mathbf{k}'} e^{-ik' \cdot y} + a^\dagger_{\mathbf{k}'} e^{ik' \cdot y} \right) \right| 0 \right\rangle.
\end{align*}
To compute the time-ordered product, we consider two cases: \( t > t' \) and \( t < t' \).
\paragraph{Case 1: \( t > t' \)}
\begin{align*}
    \langle 0 | \phi(x) \phi^\dagger(y) | 0 \rangle 
    % &= \bigints\frac{d^3\mathbf{k}}{(2\pi)^{3/2}} \frac{1}{\sqrt{2E_k}} \bigints\frac{d^3\mathbf{k}'}{(2\pi)^{3/2}} \frac{1}{\sqrt{2E_{k'}}} \left\langle 0 \left| a_{\mathbf{k}} b_{\mathbf{k}'} e^{-ik \cdot x} e^{-ik' \cdot y} + a_{\mathbf{k}} a^\dagger_{\mathbf{k}'} e^{-ik \cdot x} e^{ik' \cdot y} \right\\
    % &\qquad+ b^\dagger_{\mathbf{k}} b_{\mathbf{k}'} e^{ik \cdot x} e^{-ik' \cdot y} + b^\dagger_{\mathbf{k}} a^\dagger_{\mathbf{k}'} e^{ik \cdot x} e^{ik' \cdot y} \right| 0 \right\rangle \\
    &= \bigints\frac{d^3\mathbf{k}}{(2\pi)^{3/2}} \frac{1}{\sqrt{2E_k}} \bigints\frac{d^3\mathbf{k}'}{(2\pi)^{3/2}} \frac{1}{\sqrt{2E_{k'}}} \left\langle 0 \left| a_{\mathbf{k}} a^\dagger_{\mathbf{k}'} e^{-ik \cdot x} e^{ik' \cdot y} + b^\dagger_{\mathbf{k}} b_{\mathbf{k}'} e^{ik \cdot x} e^{-ik' \cdot y} \right| 0 \right\rangle \\
    &= \bigints\frac{d^3\mathbf{k}}{(2\pi)^{3/2}} \frac{1}{\sqrt{2E_k}} \bigints\frac{d^3\mathbf{k}'}{(2\pi)^{3/2}} \frac{1}{\sqrt{2E_{k'}}} \delta^3(\mathbf{k} - \mathbf{k'}) e^{-ik \cdot x} e^{ik \cdot y} \\
    &= \bigints\frac{d^3\mathbf{k}}{(2\pi)^{3}} \frac{1}{2E_k} e^{-ik \cdot (x - y)}.
\end{align*}

\paragraph{Case 2: \( t < t' \)}

\begin{align}
    \langle 0 | \phi^\dagger(y) \phi(x) | 0 \rangle &= \bigints\frac{d^3\mathbf{k}}{(2\pi)^{3}} \frac{1}{2E_k} e^{-ik \cdot (y - x)}.
\end{align}
Combining the results for \( t > t' \) and \( t < t' \) using time-ordering:
\begin{align}
    \langle 0 | T\{\phi(x) \phi^\dagger(y)\} | 0 \rangle &= \bigints\frac{d^3\mathbf{k}}{(2\pi)^3} \frac{1}{2E_k} \left[ \theta(t - t') e^{-iE_k (t - t')} + \theta(t' - t) e^{iE_k (t - t')} \right] e^{-i\mathbf{k} \cdot (\mathbf{x} - \mathbf{y})}.
\end{align}
In momentum space, the Feynman propagator is:
\begin{align}
    \Delta_F(p) &= \frac{1}{p^2 - m^2 + i \epsilon},
\end{align}
where \( p^2 = E^2 - \mathbf{p}^2 \) and \( \epsilon \) is a small positive number that ensures the correct analytic properties of the propagator.
\subsection*{(d) Time-Ordered Product}
To compute the time-ordered product \( T(\phi(x) \phi^\dagger(y) \phi(z)) \) for a complex scalar field \(\phi\), we need to consider all possible orderings of the fields:

\begin{align*}
T(\phi(x) \phi^\dagger(y) \phi(z)) &= \theta(t_x - t_y) \theta(t_y - t_z) \phi(x) \phi^\dagger(y) \phi(z) \\
&\quad + \theta(t_x - t_z) \theta(t_z - t_y) \phi(x) \phi(z) \phi^\dagger(y) \\
&\quad + \theta(t_y - t_x) \theta(t_x - t_z) \phi^\dagger(y) \phi(x) \phi(z) \\
&\quad + \theta(t_y - t_z) \theta(t_z - t_x) \phi^\dagger(y) \phi(z) \phi(x) \\
&\quad + \theta(t_z - t_x) \theta(t_x - t_y) \phi(z) \phi(x) \phi^\dagger(y) \\
&\quad + \theta(t_z - t_y) \theta(t_y - t_x) \phi(z) \phi^\dagger(y) \phi(x).
\end{align*}

We now evaluate each case where \( t_x \), \( t_y \), and \( t_z \) have different orderings.

\subsection*{Case 1: \( t_x > t_y > t_z \)}

\begin{align*}
\langle 0 | T(\phi(x) \phi^\dagger(y) \phi(z)) | 0 \rangle &= \langle 0 | \phi(x) \phi^\dagger(y) \phi(z) | 0 \rangle \\
&= \langle 0 | \phi(x) \phi^\dagger(y) | 0 \rangle \langle 0 | \phi(z) | 0 \rangle + \langle 0 | \phi(x) | 0 \rangle \langle 0 | \phi^\dagger(y) \phi(z) | 0 \rangle \\
&= \Delta_F(x - y) \langle 0 | \phi(z) | 0 \rangle \\
&= 0
\end{align*}

\subsection*{Case 2: \( t_x > t_z > t_y \)}

\begin{align*}
\langle 0 | T(\phi(x) \phi^\dagger(y) \phi(z)) | 0 \rangle &= \langle 0 | \phi(x) \phi(z) \phi^\dagger(y) | 0 \rangle \\
&= \langle 0 | \phi(x) | 0 \rangle \langle 0 | \phi(z) \phi^\dagger(y) | 0 \rangle + \langle 0 | \phi(x) \phi(z) | 0 \rangle \langle 0 | \phi^\dagger(y) | 0 \rangle \\
&= \Delta_F(z - y) \langle 0 | \phi(x) | 0 \rangle \\
&= 0
\end{align*}

\subsection*{Case 3: \( t_y > t_x > t_z \)}

\begin{align*}
\langle 0 | T(\phi(x) \phi^\dagger(y) \phi(z)) | 0 \rangle &= \langle 0 | \phi^\dagger(y) \phi(x) \phi(z) | 0 \rangle \\
&= \langle 0 | \phi^\dagger(y) | 0 \rangle \langle 0 | \phi(x) \phi(z) | 0 \rangle + \langle 0 | \phi^\dagger(y) \phi(x) | 0 \rangle \langle 0 | \phi(z) | 0 \rangle \\
&= \Delta_F(y - x) \langle 0 | \phi(z) | 0 \rangle \\
&= 0
\end{align*}

\subsection*{Case 4: \( t_y > t_z > t_x \)}

\begin{align*}
\langle 0 | T(\phi(x) \phi^\dagger(y) \phi(z)) | 0 \rangle &= \langle 0 | \phi^\dagger(y) \phi(z) \phi(x) | 0 \rangle \\
&= \langle 0 | \phi^\dagger(y) \phi(z) | 0 \rangle \langle 0 | \phi(x) | 0 \rangle + \langle 0 | \phi^\dagger(y) | 0 \rangle \langle 0 | \phi(z) \phi(x) | 0 \rangle \\
&= \Delta_F(z - x) \langle 0 | \phi^\dagger(y) | 0 \rangle \\
&= \Delta_F(z - x) \langle 0 | \phi^\dagger(y) | 0 \rangle
\end{align*}

\subsection*{Case 5: \( t_z > t_x > t_y \)}

\begin{align*}
\langle 0 | T(\phi(x) \phi^\dagger(y) \phi(z)) | 0 \rangle &= \langle 0 | \phi(z) \phi(x) \phi^\dagger(y) | 0 \rangle \\
&= \langle 0 | \phi(z) | 0 \rangle \langle 0 | \phi(x) \phi^\dagger(y) | 0 \rangle + \langle 0 | \phi(z) \phi(x) | 0 \rangle \langle 0 | \phi^\dagger(y) | 0 \rangle \\
&= \Delta_F(x - y) \langle 0 | \phi(z) | 0 \rangle \\
&= 0
\end{align*}

\subsection*{Case 6: \( t_z > t_y > t_x \)}

\begin{align*}
\langle 0 | T(\phi(x) \phi^\dagger(y) \phi(z)) | 0 \rangle &= \langle 0 | \phi(z) \phi^\dagger(y) \phi(x) | 0 \rangle \\
&= \langle 0 | \phi(z) \phi^\dagger(y) | 0 \rangle \langle 0 | \phi(x) | 0 \rangle + \langle 0 | \phi(z) | 0 \rangle \langle 0 | \phi^\dagger(y) \phi(x) | 0 \rangle \\
&= \Delta_F(y - x) \langle 0 | \phi(z) | 0 \rangle \\
&= \Delta_F(y - x) \langle 0 | \phi(z) | 0 \rangle
\end{align*}

In summary, the time-ordered product \( T(\phi(x) \phi^\dagger(y) \phi(z)) \) is given by:

\begin{align*}
\langle 0 | T(\phi(x) \phi^\dagger(y) \phi(z)) | 0 \rangle &= \Delta_F(x - y) \langle 0 | \phi(z) | 0 \rangle \quad \text{for } t_x > t_y > t_z, \\
\langle 0 | T(\phi(x) \phi^\dagger(y) \phi(z)) | 0 \rangle &= \Delta_F(z - y) \langle 0 | \phi(x) | 0 \rangle \quad \text{for } t_x > t_z > t_y, \\
\langle 0 | T(\phi(x) \phi^\dagger(y) \phi(z)) | 0 \rangle &= \Delta_F(y - x) \langle 0 | \phi(z) | 0 \rangle \quad \text{for } t_y > t_x > t_z, \\
\langle 0 | T(\phi(x) \phi^\dagger(y) \phi(z)) | 0 \rangle &= \Delta_F(z - x) \langle 0 | \phi^\dagger(y) | 0 \rangle \quad \text{for } t_y > t_z > t_x, \\
\langle 0 | T(\phi(x) \phi^\dagger(y) \phi(z)) | 0 \rangle &= \Delta_F(x - y) \langle 0 | \phi(z) | 0 \rangle \quad \text{for } t_z > t_x > t_y, \\
\langle 0 | T(\phi(x) \phi^\dagger(y) \phi(z)) | 0 \rangle &= \Delta_F(y - x) \langle 0 | \phi(z) | 0 \rangle \quad \text{for } t_z > t_y > t_x.
\end{align*}